\chapter{Results}
\label{ch:results}

This chapter reports what the experiments found. It is written to be
readable without prior knowledge of the method: every table and figure
is introduced in plain words before any numbers are discussed, and a set
of worked, by-hand examples (Section~\ref{sec:hand-audits}) shows what
the detector actually does on individual questions before the
aggregate scores are presented.

\section{How to read the numbers}
\label{sec:how-to-read}

The task in every experiment is the same. The language model answers a
question; sometimes it is right and sometimes it is wrong. A
\emph{detector} looks at the answer (or at several sampled answers) and
outputs a single number that is meant to be high when the answer is
unreliable. A good detector gives higher numbers to the wrong answers
than to the right ones, so that a user could decline the riskiest
answers and keep the safe ones.

Two summary measures are used throughout.

\begin{itemize}
  \item \textbf{AUROC} (area under the ROC curve). Pick one wrong answer
    and one right answer at random; the AUROC is the probability that
    the detector scored the wrong one higher. \textbf{0.5 means the
    detector is useless} (no better than a coin flip), and \textbf{1.0
    means it is perfect}. Values around 0.7--0.9 are useful in practice.
    Because AUROC only depends on the \emph{ranking} of scores, the unit
    of the entropy (nats vs.\ bits) does not affect it.
  \item \textbf{AURAC} (area under the rejection--accuracy curve). If we
    throw away the answers the detector is most worried about, the
    accuracy of what remains should go up. AURAC summarizes how quickly
    accuracy improves as we discard flagged answers; higher is better.
    Unlike AUROC, it also depends on how accurate the model was to begin
    with, so AURAC values are only comparable within the same task and
    model.
\end{itemize}

Throughout, the positive class is the \emph{incorrect} answer, so
``higher score $=$ more likely wrong''. All headline numbers are
averaged over three independent runs (random seeds 42, 43, 44); where a
spread is shown it is the standard deviation across those three runs.
The detectors compared are the ones defined in
Chapter~\ref{ch:semantic_entropy}: the cheap surface baselines
(counting distinct answer strings), discrete semantic entropy with
meaning-based clustering, the probability-weighted variants, Kernel
Language Entropy, the hidden-state probe (SEP), \mbox{P(True)}, and
embedding regression.

\section{The headline picture}
\label{sec:headline}

Figure~\ref{fig:regime-split} is the single most important result of the
thesis, and the rest of the chapter unpacks it. It compares two
detectors across the four tasks and three models: the \emph{cheap}
surface score, which just counts how many of the sampled answers are
different strings, and the \emph{meaning-aware} semantic entropy, which
first groups answers that mean the same thing and then measures the
spread over meanings.

\begin{figure}[ht]
\centering
\includegraphics[width=\textwidth]{regime_split.png}
\caption{Surface-form versus meaning-aware detection across tasks
  (AUROC, higher is better; the dashed line is the useless ``coin-flip''
  level of 0.5). On the three short-answer tasks the two scores are
  almost the same height: the cheap string-counting baseline is already
  as good as meaning-aware clustering. On the long-form biography task
  the surface bar drops to the chance line --- it stops working
  entirely --- while the meaning-aware bar stays well above it. This
  short-versus-long contrast is the central finding.}
\label{fig:regime-split}
\end{figure}

The pattern is the same for all three models. On short answers the two
bars are nearly equal, so the expensive meaning-based machinery buys
almost nothing over simply noticing that the model wrote different
strings. On long-form biographies the cheap baseline collapses to the
chance line while the meaning-aware score survives. \emph{Where you need
semantic clustering is precisely the long-form regime; on short answers
it is close to redundant.} The remainder of the chapter establishes this
quantitatively and explains why it happens.

\section{Worked examples: what the detector sees}
\label{sec:hand-audits}

Before the aggregate tables, it helps to look at three real records from
the recorded runs (canonical configuration, temperature $1.0$). Each one
shows the ten sampled answers, how the meaning-based clustering grouped
them, and the resulting scores. These by-hand audits make concrete what
the averages later summarize.

\subsection{Example 1 --- meaning-based grouping lowers a falsely high
score}
\label{sec:audit-merge}

\textbf{Question (TriviaQA, Llama-3.1-8B):} ``According to Rudyard
Kipling, what were the two impostors to meet and treat the same day?''
\textbf{Gold answer:} \emph{Triumph and Disaster}. The ten sampled
answers were:

\begin{quote}\small
\emph{Success and Failure.} $\;\cdot\;$ \emph{Triumph and Disaster}
$\;\cdot\;$ \emph{Triumph and defeat.} $\;\cdot\;$ \emph{Triumph and
disaster.} $\;\cdot\;$ \emph{Triumph and defeat.} $\;\cdot\;$
\emph{Success and Failure} $\;\cdot\;$ \emph{Success and Failure.}
$\;\cdot\;$ \emph{Triumph and Disaster} $\;\cdot\;$ \emph{Tyrant and a
dreamer.} $\;\cdot\;$ \emph{Success and Failure.}
\end{quote}

The surface baseline sees many \emph{different strings} (``Triumph and
Disaster'', ``Triumph and defeat.'', ``Success and Failure'', and so on)
and therefore reports high uncertainty: surface entropy $1.28$ nats. The
meaning-based step asks, for each pair, whether the two phrases say the
same thing in context; it judges the various ``triumph/defeat'' and
``success/failure'' phrasings to be mutually entailing and collapses
nine of the ten samples into a single meaning, leaving only the odd
``Tyrant and a dreamer'' on its own. The semantic entropy is then only
$0.33$ nats --- a much more honest reflection of the fact that the model
basically settled on one idea. The answer was graded correct. This is
the mechanism by which clustering is \emph{supposed} to help: it removes
uncertainty that is only about wording.

\subsection{Example 2 --- a confident wrong answer the method cannot
catch}
\label{sec:audit-artifact}

\textbf{Question (TriviaQA, Llama-3.1-8B):} ``Who was the director of the
CIA from 1976--81?'' \textbf{Gold answer:} \emph{George Bush}. All ten
samples were the same: \emph{Stansfield Turner}. Because every sample
agrees, the spread over meanings is zero (semantic entropy $0.00$), so
the detector reports maximum confidence --- yet the answer is graded
incorrect.

This single record illustrates two important limits at once. First, it
is a genuine \emph{confident systematic error}: the model is wrong in the
same way every time, and no consistency-based detector can flag it,
because there is no inconsistency to measure. Second, it is partly a
\emph{labelling artefact}: Stansfield Turner did head the CIA for most of
1977--1981, so the ``wrong'' label here reflects an ambiguous gold answer
as much as a model error. Both effects --- the blind spot for confident
errors and the noise in automatic correctness labels --- recur in the
aggregate numbers and are returned to in the discussion
(Chapter~\ref{ch:discussion}).

\subsection{Example 3 --- why the surface baseline dies on biographies}
\label{sec:audit-bio}

\textbf{Question (Biography, Mistral-7B):} ``Tell me a bio of Edmond
Pourchot.'' The model produced ten multi-sentence paragraphs. One began:
\emph{``Edmond Pourchot, born on March 2, 1896, in Gap, France, was a
prominent figure in the field of French aviation\ldots''} (Edmond
Pourchot was in fact a seventeenth-century French philosopher; the
paragraph is fabricated.) All ten paragraphs were different strings, so
the surface baseline reports the \emph{maximum possible} value every
time: $\ln 10 = 2.30$ nats. Across the whole biography set the average
surface entropy is $2.303$ --- pinned at the ceiling --- which is exactly
why it carries no information and lands at the chance line in
Figure~\ref{fig:regime-split}: a score that is always maximal cannot rank
anything.

Meaning-based clustering escapes this trap. It groups the paragraphs by
what they claim rather than by their exact wording, so its score varies
from record to record (average semantic entropy $1.73$, not pinned at the
ceiling). For this fabricated biography the ten paragraphs genuinely
disagree, so the semantic entropy is high and the answer is correctly
flagged as risky. This is the long-form mechanism behind the headline
figure.

\section{Short-answer results}
\label{sec:short-answer-results}

We begin with the three short-answer tasks (TriviaQA, NQ-Open, SVAMP).
Table~\ref{tab:auroc-overview} gives the bird's-eye view: discrete
semantic entropy, the full probability-weighted estimator, and the
hidden-state probe, for all three models and four tasks.

\begin{table}[ht]
\thisfloatsetup{capposition=above}
\centering
\caption{Selective-prediction AUROC overview (seed-mean over 42/43/44).
  Short-answer tasks on the left, biography on the right. On short
  answers all detectors land in a similar, useful band; on biography the
  picture changes sharply (Section~\ref{sec:bio-results}).}
\label{tab:auroc-overview}
\input{includes/generated_overview.tex}
\end{table}

\begin{figure}[ht]
\centering
\includegraphics[width=\textwidth]{multiseed_auroc.png}
\caption{Per-task AUROC for the main detectors, averaged over three
  seeds (error bars are the seed spread). The detectors cluster
  together on every short-answer task: no single method dominates, and
  the cheap surface baseline sits inside the same band as the
  meaning-aware scores.}
\label{fig:multiseed-auroc}
\end{figure}

The detailed per-task tables follow.
Tables~\ref{tab:auroc-triviaqa}, \ref{tab:auroc-nqopen}
and~\ref{tab:auroc-svamp} list every detector for TriviaQA, NQ-Open and
SVAMP respectively.

\begin{table}[ht]
\thisfloatsetup{capposition=above}
\centering
\caption{TriviaQA AUROC by detector and model (mean $\pm$ std over
  seeds). Canonical DeBERTa-v3-large clustering for the entropy methods.}
\label{tab:auroc-triviaqa}
\input{includes/generated_auroc_triviaqa.tex}
\end{table}

\begin{table}[ht]
\thisfloatsetup{capposition=above}
\centering
\caption{NQ-Open AUROC by detector and model (mean $\pm$ std over seeds).}
\label{tab:auroc-nqopen}
\input{includes/generated_auroc_nqopen.tex}
\end{table}

\begin{table}[ht]
\thisfloatsetup{capposition=above}
\centering
\caption{SVAMP AUROC by detector and model (mean $\pm$ std over seeds).}
\label{tab:auroc-svamp}
\input{includes/generated_auroc_svamp.tex}
\end{table}

Three things stand out on short answers. First, all of the
sampling-based detectors work: AUROC sits in roughly the
$0.74$--$0.88$ band, which is consistent with the original semantic
entropy paper under the matched protocol. Second, the
\emph{surface/exact-match baseline is the strongest single score in most
cells} --- for example on TriviaQA it reaches $0.819$ (Mistral),
$0.838$ (Llama-3.1-8B) and $0.802$ (Llama-3.1-70B), slightly ahead of
the meaning-aware $0.799$, $0.814$ and $0.794$. It is the single
strongest score in six of the nine short-answer cells (all of TriviaQA
and NQ-Open); only on SVAMP does the probability-weighted full estimator
edge ahead. Third, discrete semantic
entropy and the full probability-weighted estimator track each other
closely, confirming that on these tasks the cheap count-based version
loses almost nothing relative to the version that uses the model's own
probabilities. The reason, made concrete in
Example~\ref{sec:audit-merge}, is that on short answers most of the
useful grouping is already done by \emph{normalizing} the strings;
meaning-based entailment then has little left to add.

\section{Does meaning-based clustering beat the surface baseline?}
\label{sec:significance}

Because the surface and semantic scores are so close on short answers,
the honest question is whether the difference is real or just noise.
Table~\ref{tab:significance-vs-surface} reports paired bootstrap
comparisons: for each task and model it asks whether a given detector
beats the surface baseline by more than sampling noise would explain.

\begin{table}[ht]
\thisfloatsetup{capposition=above}
\centering
\caption{Paired-bootstrap comparison against the surface baseline.
  A confidence interval that contains zero means ``not significantly
  different from the cheap baseline.''}
\label{tab:significance-vs-surface}
\input{includes/generated_significance.tex}
\end{table}

On the short-answer tasks, \emph{no method beats the surface baseline by
a statistically significant margin}: wherever a method's margin over
surface entropy is large enough to be significant, it is the surface
baseline that comes out ahead, never the other way round. In other
words, once you have normalized the answer strings, paying for an NLI
model to cluster them does not reliably improve short-answer
hallucination detection.
This is a negative result, and an important one: it says the expensive
part of the method earns its keep only in the long-form setting, which
the next sections examine.

\section{Clustering backend sensitivity}
\label{sec:clustering-sensitivity}

The meaning-based step needs a model to decide whether two answers are
equivalent. Does the choice of that model matter?
Table~\ref{tab:clustering-sensitivity} compares five backends ---
exact-match, two sizes of DeBERTa NLI, and a large Qwen judge ---
holding everything else fixed.

\begin{table}[ht]
\thisfloatsetup{capposition=above}
\centering
\caption{Discrete semantic entropy AUROC by clustering backend
  (mean $\pm$ std over seeds). ``--'' marks a backend that was not run
  for that task (the larger NLI judges were applied to short answers
  only as a sensitivity check).}
\label{tab:clustering-sensitivity}
\input{includes/generated_clustering_sensitivity.tex}
\end{table}

\begin{figure}[ht]
\centering
\includegraphics[width=0.85\textwidth]{bio_clustering_sensitivity.png}
\caption{Biography AUROC for discrete semantic entropy under different
  clustering backends. On long-form inputs the backend is
  result-determining: the same samples yield anything from chance-level
  to clearly useful detection depending only on how equivalence is
  judged.}
\label{fig:bio-clustering-sensitivity}
\end{figure}

On short answers the backend barely matters --- exact-match, base NLI,
large NLI and the Qwen judge land within a few points of each other,
which is unsurprising given that normalization already does the work. On
biographies the choice is decisive: the Qwen judge clusters meanings more
reliably than the smaller DeBERTa models and lifts AUROC accordingly.
This is the long-form counterpart to the short-answer negative result:
when clustering matters at all, \emph{which} clusterer you use matters a
great deal.

\section{Long-form biography results}
\label{sec:bio-results}

The biography task is the controlled long-form case: the same
sample-and-cluster pipeline, but the model now writes a paragraph rather
than a phrase. Table~\ref{tab:auroc-bio} lists every detector for all
three generators.

\begin{table}[ht]
\thisfloatsetup{capposition=above}
\centering
\caption{Biography (whole-paragraph) AUROC by detector and model
  (mean $\pm$ std over seeds), canonical DeBERTa-v3-large clustering.
  The two count-based surface scores (naive \emph{sample} entropy and
  surface entropy) collapse to chance ($\approx 0.50$) because long
  paragraphs almost never repeat as identical strings; the
  probability-weighted naive entropy, the meaning-aware methods, and the
  hidden-state probe all recover a real signal.}
\label{tab:auroc-bio}
\input{includes/generated_auroc_bio.tex}
\end{table}

The contrast with short answers is stark. The two count-based surface
scores --- naive \emph{sample} entropy and surface entropy --- sit at
$\approx 0.50$, pure chance, exactly as Example~\ref{sec:audit-bio}
predicted: every long paragraph is a distinct string, so these scores
are permanently pinned at their maximum and can rank nothing. Every
detector that escapes raw string-counting recovers a real signal.
Meaning-aware discrete semantic entropy reaches $0.708$ (Mistral-7B),
$0.618$ (Llama-3.1-8B) and $0.612$ (Llama-3.1-70B), and the Qwen-judge
clustering variant is stronger still ($0.736 / 0.685 / 0.666$),
consistent with the backend sensitivity above. So the value of moving
from surface comparison to meaning comparison is not marginal here --- it
is the difference between a useless detector and a working one.

Two detectors that need no clustering at all are in fact the strongest on
biography. Probability-weighted naive entropy --- which keeps each
distinct paragraph but weights it by the model's own sequence
probability --- is the single best score on the two smaller models
($0.781$ Mistral, $0.769$ Llama-3.1-8B), and the self-evaluation
baseline \mbox{P(True)} is the best on the largest ($0.777$). The lesson
is that on long text the model's \emph{own probability signal} carries
much of the confabulation information, and the hard-clustering step only
partially recovers it.

One regularity is worth flagging. The semantic-entropy AUROC
\emph{decreases} as the generator gets larger ($0.708 \to 0.618 \to
0.612$ from 7B to 70B). The discussion returns to this, but the short
explanation is that larger models confabulate more \emph{consistently}:
when a big model invents a fact, its ten samples tend to agree on the
invention, leaving little spread for a consistency-based detector to
find.

\section{Supervised baselines and hidden-state probes}
\label{sec:probes-results}

Semantic entropy needs several samples per question. Two cheaper-at-test
alternatives read the model's internal state from a single pass: the
hidden-state probe (SEP) and embedding regression.
Table~\ref{tab:embedding-regression} reports them alongside
\mbox{P(True)}.

\begin{table}[ht]
\thisfloatsetup{capposition=above}
\centering
\caption{Probe and supervised-baseline AUROC (mean $\pm$ std over
  seeds): hidden-state SEP probe, embedding regression
  (in-distribution and out-of-distribution), and \mbox{P(True)}.}
\label{tab:embedding-regression}
\input{includes/generated_embedding_regression.tex}
\end{table}

The probes are competitive, and on biography they are strong: the SEP
probe reaches $0.728$ (8B) and $0.731$ (70B), \emph{above} every
meaning-based clustering method on those models, and \mbox{P(True)} is
the single best biography detector on the 70B model at $0.777$. (Only
the probability-weighted naive entropy, noted above, beats the probe on
the two smaller models.) This matters for deployment, because a
single-pass probe is far cheaper than drawing ten samples and running
pairwise entailment. The catch, shown next, is that the probe's
advantage depends on the decoding temperature.

\section{Sampling-temperature sensitivity}
\label{sec:temperature-sensitivity}

Every result so far used the model's default sampling temperature. Since
sampling-based detectors depend on getting \emph{diverse} samples, the
decoding temperature ought to matter --- and it does, in a way that has a
clean practical lesson. Table~\ref{tab:temp-sweep} and
Figure~\ref{fig:temp-sweep} sweep the temperature from $0.1$ (nearly
deterministic) to $1.0$.

\begin{table}[ht]
\thisfloatsetup{capposition=above}
\centering
\caption{AUROC versus sampling temperature (seed-mean), Llama-3.1-70B.
  Sampling-based entropy rises with temperature; the single-pass probe
  falls.}
\label{tab:temp-sweep}
\input{includes/generated_temp_sweep_auroc.tex}
\end{table}

\begin{figure}[ht]
\centering
\includegraphics[width=\textwidth]{auroc_vs_temperature.png}
\caption{AUROC as a function of sampling temperature, per task and
  model. The sampling-based scores (semantic entropy, surface entropy)
  climb as temperature increases, because more diverse samples reveal
  more disagreement to measure. The hidden-state probe moves the
  opposite way. The two families therefore \emph{cross over} between
  temperature $0.7$ and $1.0$ on short answers; on biography the curves
  are flat.}
\label{fig:temp-sweep}
\end{figure}

At low temperature the model repeats nearly the same answer every time,
so there is little diversity for a sampling-based method to work with and
semantic entropy is weak (TriviaQA/70B: AUROC $0.54$ at $T{=}0.1$). As
temperature rises the samples spread out and the same method becomes
strong ($0.78$ at $T{=}1.0$). The hidden-state probe does the reverse ---
it is best at low temperature and degrades as decoding gets noisier ---
so the two families \emph{cross over}. The practical lesson is that there
is no single best detector: at low decoding temperature, use the probe;
at high temperature, use sampling-based entropy. On biography the curves
are flat, so this trade-off is specific to the short-answer regime.

\section{Claim-level long-form detection: FactualBio}
\label{sec:factualbio-results}

The whole-paragraph biography score asks ``is this paragraph reliable?''
A finer question is ``which individual claims in the paragraph are
reliable?'' The FactualBio extension breaks each biography into
atomic claims and applies the same meaning-based idea at the claim level.
Table~\ref{tab:factualbio} reports the result.

\begin{table}[ht]
\thisfloatsetup{capposition=above}
\centering
\caption{Claim-level FactualBio AUROC (NLI clustering), by generator.
  Detection quality decreases as the generator gets larger.}
\label{tab:factualbio}
\input{includes/generated_factualbio_main.tex}
\end{table}

At the claim level the same scale effect from
Section~\ref{sec:bio-results} appears even more cleanly: claim-level
detection AUROC falls monotonically with model size, $0.688$ (Mistral-7B)
$\to 0.644$ (Llama-3.1-8B) $\to 0.521$ (Llama-3.1-70B). By the time we
reach the 70B model, claim-level semantic entropy is barely above chance.
Larger models repeat their confabulations too consistently for a
consistency-based signal to separate true from false claims.

\section{Summary of findings}
\label{sec:results-summary}

The experiments support a single, regime-dependent picture.

\begin{enumerate}
  \item \textbf{Short answers: semantic clustering is close to
    redundant.} All sampling-based detectors work (AUROC
    $\approx 0.74$--$0.88$), but the cheap surface/exact-match baseline
    is the strongest single score in most cells and \emph{no method
    beats it significantly} (Section~\ref{sec:significance}).
    Normalization, not entailment, does most of the grouping.
  \item \textbf{Long-form: semantic clustering is decisive.} Surface
    baselines collapse to chance on biographies, while meaning-aware
    entropy, the LLM-judge variant, and the hidden-state probe all
    recover a real signal (Section~\ref{sec:bio-results}). The clustering
    backend is result-determining here
    (Section~\ref{sec:clustering-sensitivity}).
  \item \textbf{Scale erodes the signal.} Both whole-paragraph and
    claim-level detection get \emph{worse} as the generator grows,
    because larger models confabulate more consistently
    (Section~\ref{sec:factualbio-results}).
  \item \textbf{Temperature decides which detector wins.} Sampling-based
    entropy and the single-pass probe respond to decoding temperature in
    opposite directions and cross over on short answers
    (Section~\ref{sec:temperature-sensitivity}).
  \item \textbf{Two blind spots, seen by hand.} Confident systematic
    errors and noisy gold labels (Example~\ref{sec:audit-artifact}) are
    real limits that the aggregate scores hide; they are taken up in
    Chapter~\ref{ch:discussion}.
\end{enumerate}
